\documentclass[11pt,a4paper]{article}
\usepackage[utf8]{inputenc}
\usepackage[T1]{fontenc}
\usepackage[french]{babel}
\usepackage{amsmath}
\usepackage{amssymb}
\usepackage{amsthm}
\usepackage{mathtools}
\usepackage{geometry}
\usepackage{enumitem}
\usepackage{xcolor}
\usepackage{titlesec}
\usepackage{fancyhdr}
\usepackage{booktabs}
\usepackage{array}
\usepackage{longtable}
\usepackage{tcolorbox}
\usepackage{multicol}

% Configuration de la page
\geometry{margin=2.5cm}
\pagestyle{fancy}
\fancyhf{}
\fancyhead[L]{\leftmark}
\fancyhead[R]{Exercices Tests d'Hypothèses}
\fancyfoot[C]{\thepage}

% Configuration des titres
\titleformat{\section}
{\Large\bfseries\color{blue!70!black}}
{}
{0em}
{}[\titlerule]

\titleformat{\subsection}
{\large\bfseries\color{blue!50!black}}
{}
{0em}
{}

\titleformat{\subsubsection}
{\normalsize\bfseries\color{blue!30!black}}
{}
{0em}
{}

% Boîtes colorées pour astuces et rappels
\newtcolorbox{astucebox}{
    colback=green!5!white,
    colframe=green!75!black,
    title=\textbf{Astuce},
    fonttitle=\bfseries
}

\newtcolorbox{rappelbox}{
    colback=blue!5!white,
    colframe=blue!75!black,
    title=\textbf{Rappel},
    fonttitle=\bfseries
}

\newtcolorbox{attentionbox}{
    colback=orange!5!white,
    colframe=orange!75!black,
    title=\textbf{Attention},
    fonttitle=\bfseries
}

% Commandes personnalisées
\newcommand{\R}{\mathbb{R}}
\newcommand{\Var}{\text{Var}}
\newcommand{\E}{\mathbb{E}}
\newcommand{\N}{\mathcal{N}}
\newcommand{\chideux}{\chi^2}
\newcommand{\student}{t}

% Métadonnées
\title{Exercices Complets - Tests d'Hypothèses\\
\large Khi-deux, P-value, Moyenne d'un Échantillon}
\author{AmineGR03}
\date{\today}

\begin{document}

\maketitle

\tableofcontents
\newpage

\section{Introduction et Méthodologie}

\subsection{Objectif de ce document}

Ce document présente une série d'exercices pratiques sur trois concepts fondamentaux des tests d'hypothèses :
\begin{itemize}
    \item \textbf{Test du Khi-deux} ($\chideux$) : pour tester la variance ou l'adéquation à une distribution
    \item \textbf{P-value} : probabilité d'observer des résultats aussi extrêmes sous $H_0$
    \item \textbf{Test sur la moyenne} : tests unilatéraux et bilatéraux
\end{itemize}

\subsection{Méthodologie générale}

Pour résoudre un problème de test d'hypothèse, suivez ces étapes :

\begin{enumerate}
    \item \textbf{Identifier le paramètre à tester} (moyenne $\mu$, variance $\sigma^2$, etc.)
    \item \textbf{Formuler les hypothèses} $H_0$ et $H_1$
    \item \textbf{Choisir la statistique de test} appropriée
    \item \textbf{Déterminer la distribution} sous $H_0$ (Normale, Student-t, $\chideux$)
    \item \textbf{Calculer la valeur observée} de la statistique
    \item \textbf{Utiliser les tables} ou calculer la p-value
    \item \textbf{Prendre une décision} (rejeter ou accepter $H_0$)
    \item \textbf{Interpréter les résultats} dans le contexte
\end{enumerate}

\newpage

\section{Partie 1 : Test du Khi-deux ($\chideux$)}

\subsection{Théorie et Étapes à Suivre}

\subsubsection{Définition}

Le test du Khi-deux est utilisé pour :
\begin{itemize}
    \item Tester la variance d'une population normale : $H_0 : \sigma^2 = \sigma_0^2$
    \item Tester l'adéquation à une distribution (test d'ajustement)
    \item Tester l'indépendance entre variables (tableau de contingence)
\end{itemize}

\subsubsection{Étapes pour tester la variance}

\begin{enumerate}
    \item \textbf{Formuler les hypothèses} :
    \begin{align*}
    H_0 &: \sigma^2 = \sigma_0^2 \\
    H_1 &: \sigma^2 \neq \sigma_0^2 \quad \text{(bilatéral)} \\
    \text{ou} \quad H_1 &: \sigma^2 < \sigma_0^2 \quad \text{(unilatéral à gauche)} \\
    \text{ou} \quad H_1 &: \sigma^2 > \sigma_0^2 \quad \text{(unilatéral à droite)}
    \end{align*}
    
    \item \textbf{Statistique de test} :
    \begin{align*}
    \chideux_{\text{obs}} = \frac{(n-1)S^2}{\sigma_0^2}
    \end{align*}
    où $S^2$ est la variance de l'échantillon et $n$ la taille de l'échantillon.
    
    \item \textbf{Distribution sous $H_0$} :
    \begin{align*}
    \chideux_{\text{obs}} \sim \chideux(n-1)
    \end{align*}
    Suit une loi du Khi-deux à $(n-1)$ degrés de liberté.
    
    \item \textbf{Utilisation de la table du Khi-deux} :
    \begin{itemize}
        \item Pour un test \textbf{bilatéral} au seuil $\alpha$ :
        \begin{align*}
        \text{Rejeter } H_0 \text{ si } \chideux_{\text{obs}} < \chideux_{1-\alpha/2}(n-1) \text{ ou } \chideux_{\text{obs}} > \chideux_{\alpha/2}(n-1)
        \end{align*}
        \item Pour un test \textbf{unilatéral à gauche} :
        \begin{align*}
        \text{Rejeter } H_0 \text{ si } \chideux_{\text{obs}} < \chideux_{1-\alpha}(n-1)
        \end{align*}
        \item Pour un test \textbf{unilatéral à droite} :
        \begin{align*}
        \text{Rejeter } H_0 \text{ si } \chideux_{\text{obs}} > \chideux_{\alpha}(n-1)
        \end{align*}
    \end{itemize}
    
    \item \textbf{Calcul de la p-value} :
    \begin{itemize}
        \item Bilatéral : $p = 2 \times \min(P(\chideux < \chideux_{\text{obs}}), P(\chideux > \chideux_{\text{obs}}))$
        \item Unilatéral : $p = P(\chideux > \chideux_{\text{obs}})$ ou $p = P(\chideux < \chideux_{\text{obs}})$
    \end{itemize}
\end{enumerate}

\begin{astucebox}
\textbf{Astuce pour lire la table du Khi-deux} :
\begin{itemize}
    \item La table donne les valeurs critiques $\chideux_{\alpha}(\nu)$ où $\nu$ est le nombre de degrés de liberté
    \item $\chideux_{\alpha}(\nu)$ est la valeur telle que $P(\chideux > \chideux_{\alpha}(\nu)) = \alpha$
    \item Pour $\chideux_{1-\alpha}(\nu)$, cherchez dans la colonne correspondant à $(1-\alpha)$
    \item Si $\nu > 30$, on peut approximer par une loi normale : $\chideux(\nu) \approx \N(\nu, 2\nu)$
\end{itemize}
\end{astucebox}

\newpage

\subsection{Exercices sur le Khi-deux}

\subsubsection{Exercice 1 : Test bilatéral sur la variance}

Une machine produit des pièces dont le diamètre suit une loi normale. Le fabricant prétend que l'écart-type est de 2 mm. On prélève un échantillon de 25 pièces et on trouve un écart-type de 2.5 mm.

\textbf{Question} : Peut-on rejeter l'hypothèse que l'écart-type est de 2 mm au seuil de signification de 5\% ?

\subsubsection*{Solution de l'Exercice 1}

\begin{enumerate}
    \item \textbf{Formulation des hypothèses} :
    \begin{align*}
    H_0 &: \sigma^2 = 4 \quad \text{(car } \sigma = 2 \Rightarrow \sigma^2 = 4\text{)} \\
    H_1 &: \sigma^2 \neq 4 \quad \text{(test bilatéral)}
    \end{align*}
    
    \item \textbf{Données} :
    \begin{itemize}
        \item $n = 25$ (taille de l'échantillon)
        \item $S = 2.5$ mm (écart-type observé)
        \item $S^2 = 6.25$ mm²
        \item $\sigma_0^2 = 4$ mm²
        \item $\alpha = 0.05$ (seuil de signification)
    \end{itemize}
    
    \item \textbf{Calcul de la statistique de test} :
    \begin{align*}
    \chideux_{\text{obs}} &= \frac{(n-1)S^2}{\sigma_0^2} \\
    &= \frac{(25-1) \times 6.25}{4} \\
    &= \frac{24 \times 6.25}{4} \\
    &= \frac{150}{4} = 37.5
    \end{align*}
    
    \item \textbf{Degrés de liberté} : $\nu = n - 1 = 25 - 1 = 24$
    
    \item \textbf{Valeurs critiques} (test bilatéral, $\alpha = 0.05$) :
    \begin{align*}
    \chideux_{0.975}(24) &= 12.401 \quad \text{(borne inférieure)} \\
    \chideux_{0.025}(24) &= 39.364 \quad \text{(borne supérieure)}
    \end{align*}
    
    \item \textbf{Décision} :
    \begin{align*}
    \chideux_{\text{obs}} = 37.5
    \end{align*}
    On a : $12.401 < 37.5 < 39.364$
    
    Donc on \textbf{accepte} $H_0$ au seuil de 5\%. L'écart-type observé n'est pas significativement différent de 2 mm.
    
    \item \textbf{Interprétation} : Les données ne permettent pas de rejeter l'hypothèse que l'écart-type est de 2 mm. La différence observée peut être due au hasard de l'échantillonnage.
\end{enumerate}

\begin{rappelbox}
\textbf{Rappel} : Pour un test bilatéral, on rejette $H_0$ si la statistique observée est en dehors de l'intervalle $[\chideux_{1-\alpha/2}, \chideux_{\alpha/2}]$.
\end{rappelbox}

\newpage

\subsubsection{Exercice 2 : Test unilatéral sur la variance}

Un contrôleur qualité veut vérifier si la variance de la production a diminué après une amélioration du processus. Avant l'amélioration, la variance était de 10. On prélève un échantillon de 20 unités et on trouve une variance de 6.5.

\textbf{Question} : Peut-on conclure que la variance a diminué au seuil de 1\% ?

\subsubsection*{Solution de l'Exercice 2}

\begin{enumerate}
    \item \textbf{Formulation des hypothèses} :
    \begin{align*}
    H_0 &: \sigma^2 = 10 \\
    H_1 &: \sigma^2 < 10 \quad \text{(test unilatéral à gauche)}
    \end{align*}
    
    \item \textbf{Données} :
    \begin{itemize}
        \item $n = 20$
        \item $S^2 = 6.5$
        \item $\sigma_0^2 = 10$
        \item $\alpha = 0.01$
    \end{itemize}
    
    \item \textbf{Calcul de la statistique de test} :
    \begin{align*}
    \chideux_{\text{obs}} &= \frac{(n-1)S^2}{\sigma_0^2} \\
    &= \frac{(20-1) \times 6.5}{10} \\
    &= \frac{19 \times 6.5}{10} \\
    &= \frac{123.5}{10} = 12.35
    \end{align*}
    
    \item \textbf{Degrés de liberté} : $\nu = 20 - 1 = 19$
    
    \item \textbf{Valeur critique} (test unilatéral à gauche, $\alpha = 0.01$) :
    \begin{align*}
    \chideux_{0.99}(19) = 7.633
    \end{align*}
    
    \item \textbf{Décision} :
    \begin{align*}
    \chideux_{\text{obs}} = 12.35 > 7.633 = \chideux_{0.99}(19)
    \end{align*}
    
    On \textbf{accepte} $H_0$ au seuil de 1\%. On ne peut pas conclure que la variance a diminué.
    
    \item \textbf{Interprétation} : Les données ne montrent pas de diminution significative de la variance. L'amélioration n'a peut-être pas eu l'effet escompté, ou l'échantillon est trop petit pour détecter la différence.
\end{enumerate}

\begin{attentionbox}
\textbf{Attention} : Pour un test unilatéral à gauche ($H_1 : \sigma^2 < \sigma_0^2$), on rejette $H_0$ si $\chideux_{\text{obs}} < \chideux_{1-\alpha}(\nu)$. Une petite valeur de $\chideux_{\text{obs}}$ correspond à une petite variance observée.
\end{attentionbox}

\newpage

\section{Partie 2 : P-value}

\subsection{Théorie et Étapes à Suivre}

\subsubsection{Définition}

La \textbf{p-value} (probabilité critique) est la probabilité d'observer une valeur de la statistique de test aussi extrême, ou plus extrême, que celle observée, sous l'hypothèse que $H_0$ est vraie.

\subsubsection{Interprétation}

\begin{itemize}
    \item \textbf{p-value < $\alpha$} : On rejette $H_0$ (résultat significatif)
    \item \textbf{p-value $\geq \alpha$} : On accepte $H_0$ (résultat non significatif)
    \item Plus la p-value est petite, plus la preuve contre $H_0$ est forte
\end{itemize}

\subsubsection{Étapes pour calculer la p-value}

\begin{enumerate}
    \item \textbf{Calculer la statistique de test observée}
    \item \textbf{Identifier la distribution} sous $H_0$
    \item \textbf{Déterminer le type de test} (unilatéral ou bilatéral)
    \item \textbf{Calculer la probabilité} :
    \begin{itemize}
        \item Test \textbf{bilatéral} : $p = 2 \times P(\text{statistique} > |\text{valeur observée}|)$
        \item Test \textbf{unilatéral à droite} : $p = P(\text{statistique} > \text{valeur observée})$
        \item Test \textbf{unilatéral à gauche} : $p = P(\text{statistique} < \text{valeur observée})$
    \end{itemize}
    \item \textbf{Utiliser les tables} ou un logiciel pour obtenir la probabilité
    \item \textbf{Comparer avec $\alpha$} et prendre une décision
\end{enumerate}

\begin{astucebox}
\textbf{Astuce pour calculer la p-value avec les tables} :
\begin{itemize}
    \item Pour la loi normale : utilisez la table de la fonction de répartition $\Phi(z)$
    \item Pour la loi de Student : la table donne souvent $P(t > t_{\alpha})$ directement
    \item Pour le Khi-deux : la table donne $P(\chideux > \chideux_{\alpha})$
    \item Si la valeur observée est entre deux valeurs de la table, faites une interpolation linéaire
    \item Pour les tests bilatéraux, n'oubliez pas de multiplier par 2 !
\end{itemize}
\end{astucebox}

\newpage

\subsection{Exercices sur la P-value}

\subsubsection{Exercice 3 : P-value pour un test sur la moyenne (loi normale)}

On teste $H_0 : \mu = 100$ contre $H_1 : \mu \neq 100$ avec un échantillon de taille $n = 36$ d'une population normale de variance connue $\sigma^2 = 25$. La moyenne observée est $\bar{x} = 102.5$.

\textbf{Question} : Calculer la p-value et conclure au seuil de 5\%.

\subsubsection*{Solution de l'Exercice 3}

\begin{enumerate}
    \item \textbf{Statistique de test} (loi normale, variance connue) :
    \begin{align*}
    Z_{\text{obs}} &= \frac{\bar{X} - \mu_0}{\sigma/\sqrt{n}} \\
    &= \frac{102.5 - 100}{5/\sqrt{36}} \\
    &= \frac{2.5}{5/6} \\
    &= \frac{2.5 \times 6}{5} = \frac{15}{5} = 3
    \end{align*}
    
    \item \textbf{Distribution sous $H_0$} : $Z_{\text{obs}} \sim \N(0,1)$
    
    \item \textbf{Calcul de la p-value} (test bilatéral) :
    \begin{align*}
    p &= 2 \times P(Z > |Z_{\text{obs}}|) \\
    &= 2 \times P(Z > 3) \\
    &= 2 \times (1 - \Phi(3)) \\
    &= 2 \times (1 - 0.9987) \\
    &= 2 \times 0.0013 = 0.0026
    \end{align*}
    
    \item \textbf{Décision} :
    \begin{align*}
    p = 0.0026 < 0.05 = \alpha
    \end{align*}
    
    On \textbf{rejette} $H_0$ au seuil de 5\%. La moyenne est significativement différente de 100.
    
    \item \textbf{Interprétation} : La probabilité d'observer une moyenne aussi éloignée de 100 sous $H_0$ est très faible (0.26\%). Cela constitue une preuve forte contre $H_0$.
\end{enumerate}

\begin{rappelbox}
\textbf{Rappel} : Pour lire $\Phi(3)$ dans la table normale :
\begin{itemize}
    \item Cherchez la ligne correspondant à $z = 3.0$
    \item La valeur est environ 0.9987
    \item $P(Z > 3) = 1 - \Phi(3) = 1 - 0.9987 = 0.0013$
\end{itemize}
\end{rappelbox}

\newpage

\subsubsection{Exercice 4 : P-value pour un test unilatéral (loi de Student)}

On teste $H_0 : \mu = 50$ contre $H_1 : \mu > 50$ avec un échantillon de taille $n = 16$ d'une population normale de variance inconnue. La moyenne observée est $\bar{x} = 52.8$ et l'écart-type observé est $s = 4.2$.

\textbf{Question} : Calculer la p-value et conclure au seuil de 1\%.

\subsubsection*{Solution de l'Exercice 4}

\begin{enumerate}
    \item \textbf{Statistique de test} (loi de Student, variance inconnue) :
    \begin{align*}
    t_{\text{obs}} &= \frac{\bar{X} - \mu_0}{S/\sqrt{n}} \\
    &= \frac{52.8 - 50}{4.2/\sqrt{16}} \\
    &= \frac{2.8}{4.2/4} \\
    &= \frac{2.8}{1.05} = 2.667
    \end{align*}
    
    \item \textbf{Degrés de liberté} : $\nu = n - 1 = 16 - 1 = 15$
    
    \item \textbf{Distribution sous $H_0$} : $t_{\text{obs}} \sim t(15)$
    
    \item \textbf{Calcul de la p-value} (test unilatéral à droite) :
    \begin{align*}
    p &= P(t(15) > t_{\text{obs}}) \\
    &= P(t(15) > 2.667)
    \end{align*}
    
    Dans la table de Student avec $\nu = 15$ :
    \begin{itemize}
        \item $t_{0.01}(15) = 2.602$
        \item $t_{0.005}(15) = 2.947$
    \end{itemize}
    
    On a $2.602 < 2.667 < 2.947$, donc :
    \begin{align*}
    0.005 < p < 0.01
    \end{align*}
    
    Par interpolation linéaire :
    \begin{align*}
    p &\approx 0.005 + \frac{2.667 - 2.602}{2.947 - 2.602} \times (0.01 - 0.005) \\
    &\approx 0.005 + \frac{0.065}{0.345} \times 0.005 \\
    &\approx 0.005 + 0.00094 \approx 0.00594
    \end{align*}
    
    \item \textbf{Décision} :
    \begin{align*}
    p \approx 0.00594 < 0.01 = \alpha
    \end{align*}
    
    On \textbf{rejette} $H_0$ au seuil de 1\%. La moyenne est significativement supérieure à 50.
    
    \item \textbf{Interprétation} : La probabilité d'observer une moyenne aussi élevée sous $H_0$ est inférieure à 1\%. C'est une preuve forte que $\mu > 50$.
\end{enumerate}

\begin{astucebox}
\textbf{Astuce pour l'interpolation} : Si $t_{\text{obs}}$ est entre $t_{\alpha_1}$ et $t_{\alpha_2}$ dans la table, la p-value est approximativement entre $\alpha_1$ et $\alpha_2$. Pour une meilleure précision, utilisez l'interpolation linéaire.
\end{astucebox}

\newpage

\section{Partie 3 : Test sur la Moyenne d'un Échantillon}

\subsection{Théorie et Étapes à Suivre}

\subsubsection{Choix de la distribution}

Le choix de la distribution dépend de ce que nous connaissons sur la population :

\begin{enumerate}
    \item \textbf{Variance connue} $\Rightarrow$ Utiliser la \textbf{loi normale} $\N(0,1)$
    \item \textbf{Variance inconnue} $\Rightarrow$ Utiliser la \textbf{loi de Student-t}
    \item \textbf{Grand échantillon} ($n \geq 30$) $\Rightarrow$ On peut approximer par la loi normale même si la variance est inconnue
\end{enumerate}

\subsubsection{Statistiques de test}

\begin{itemize}
    \item \textbf{Variance connue} :
    \begin{align*}
    Z = \frac{\bar{X} - \mu_0}{\sigma/\sqrt{n}} \sim \N(0,1)
    \end{align*}
    
    \item \textbf{Variance inconnue} :
    \begin{align*}
    T = \frac{\bar{X} - \mu_0}{S/\sqrt{n}} \sim t(n-1)
    \end{align*}
    où $S$ est l'écart-type de l'échantillon.
\end{itemize}

\subsubsection{Étapes pour un test sur la moyenne}

\begin{enumerate}
    \item \textbf{Identifier les informations disponibles} :
    \begin{itemize}
        \item Variance connue ou inconnue ?
        \item Taille de l'échantillon ?
        \item Distribution de la population (normale ou non) ?
    \end{itemize}
    
    \item \textbf{Choisir la distribution appropriée} :
    \begin{itemize}
        \item Variance connue $\Rightarrow$ Table normale
        \item Variance inconnue + $n < 30$ $\Rightarrow$ Table de Student
        \item Variance inconnue + $n \geq 30$ $\Rightarrow$ Table normale (approximation)
    \end{itemize}
    
    \item \textbf{Formuler les hypothèses} :
    \begin{itemize}
        \item Bilatéral : $H_0 : \mu = \mu_0$ vs $H_1 : \mu \neq \mu_0$
        \item Unilatéral à gauche : $H_0 : \mu = \mu_0$ vs $H_1 : \mu < \mu_0$
        \item Unilatéral à droite : $H_0 : \mu = \mu_0$ vs $H_1 : \mu > \mu_0$
    \end{itemize}
    
    \item \textbf{Calculer la statistique de test}
    
    \item \textbf{Utiliser la table appropriée} pour trouver la valeur critique ou la p-value
    
    \item \textbf{Prendre une décision}
    
    \item \textbf{Interpréter les résultats}
\end{enumerate}

\begin{rappelbox}
\textbf{Rappel sur les tables} :
\begin{itemize}
    \item \textbf{Table normale} : Donne $\Phi(z) = P(Z \leq z)$
    \item \textbf{Table de Student} : Donne $t_{\alpha}(\nu)$ tel que $P(t(\nu) > t_{\alpha}(\nu)) = \alpha$
    \item Pour un test bilatéral, utilisez $\alpha/2$ dans la table
    \item Pour un test unilatéral, utilisez $\alpha$ directement
\end{itemize}
\end{rappelbox}

\newpage

\subsection{Exercices : Test Bilatéral sur la Moyenne}

\subsubsection{Exercice 5 : Test bilatéral avec variance connue}

Un fabricant de batteries prétend que ses batteries ont une durée de vie moyenne de 120 heures. On teste un échantillon de 50 batteries et on trouve une moyenne de 118 heures. On suppose que la durée de vie suit une loi normale avec un écart-type de 10 heures.

\textbf{Question} : Peut-on rejeter l'affirmation du fabricant au seuil de 5\% ?

\subsubsection*{Solution de l'Exercice 5}

\begin{enumerate}
    \item \textbf{Formulation des hypothèses} :
    \begin{align*}
    H_0 &: \mu = 120 \\
    H_1 &: \mu \neq 120 \quad \text{(test bilatéral)}
    \end{align*}
    
    \item \textbf{Données} :
    \begin{itemize}
        \item $n = 50$ (grand échantillon)
        \item $\bar{x} = 118$ heures
        \item $\sigma = 10$ heures (connu)
        \item $\mu_0 = 120$ heures
        \item $\alpha = 0.05$
    \end{itemize}
    
    \item \textbf{Choix de la distribution} : Variance connue $\Rightarrow$ Loi normale $\N(0,1)$
    
    \item \textbf{Calcul de la statistique de test} :
    \begin{align*}
    Z_{\text{obs}} &= \frac{\bar{X} - \mu_0}{\sigma/\sqrt{n}} \\
    &= \frac{118 - 120}{10/\sqrt{50}} \\
    &= \frac{-2}{10/7.071} \\
    &= \frac{-2}{1.414} = -1.414
    \end{align*}
    
    \item \textbf{Valeurs critiques} (test bilatéral, $\alpha = 0.05$) :
    \begin{align*}
    z_{\alpha/2} = z_{0.025} = 1.96
    \end{align*}
    
    Dans la table normale : $\Phi(1.96) = 0.975$, donc $P(Z > 1.96) = 0.025$
    
    \item \textbf{Décision} :
    \begin{align*}
    |Z_{\text{obs}}| = 1.414 < 1.96 = z_{0.025}
    \end{align*}
    
    On \textbf{accepte} $H_0$ au seuil de 5\%. On ne peut pas rejeter l'affirmation du fabricant.
    
    \item \textbf{Calcul de la p-value} :
    \begin{align*}
    p &= 2 \times P(Z > |Z_{\text{obs}}|) \\
    &= 2 \times P(Z > 1.414) \\
    &= 2 \times (1 - \Phi(1.414)) \\
    &= 2 \times (1 - 0.9213) \\
    &= 2 \times 0.0787 = 0.1574
    \end{align*}
    
    $p = 0.1574 > 0.05 = \alpha$, ce qui confirme qu'on accepte $H_0$.
    
    \item \textbf{Interprétation} : La différence observée (118 vs 120) n'est pas statistiquement significative. Elle peut être due au hasard de l'échantillonnage.
\end{enumerate}

\begin{astucebox}
\textbf{Astuce} : Pour un test bilatéral avec la loi normale, on compare toujours la valeur absolue de la statistique observée avec la valeur critique. La règle est : rejeter $H_0$ si $|Z_{\text{obs}}| > z_{\alpha/2}$.
\end{astucebox}

\newpage

\subsubsection{Exercice 6 : Test bilatéral avec variance inconnue (loi de Student)}

Un laboratoire pharmaceutique prétend qu'un médicament réduit la tension artérielle de 10 mmHg en moyenne. On teste le médicament sur 20 patients et on obtient une réduction moyenne de 8.5 mmHg avec un écart-type de 3.2 mmHg. On suppose que les réductions suivent une loi normale.

\textbf{Question} : Les données confirment-elles l'affirmation du laboratoire au seuil de 5\% ?

\subsubsection*{Solution de l'Exercice 6}

\begin{enumerate}
    \item \textbf{Formulation des hypothèses} :
    \begin{align*}
    H_0 &: \mu = 10 \\
    H_1 &: \mu \neq 10 \quad \text{(test bilatéral)}
    \end{align*}
    
    \item \textbf{Données} :
    \begin{itemize}
        \item $n = 20$ (petit échantillon)
        \item $\bar{x} = 8.5$ mmHg
        \item $s = 3.2$ mmHg
        \item $\mu_0 = 10$ mmHg
        \item $\alpha = 0.05$
    \end{itemize}
    
    \item \textbf{Choix de la distribution} : Variance inconnue + $n < 30$ $\Rightarrow$ Loi de Student $t(19)$
    
    \item \textbf{Calcul de la statistique de test} :
    \begin{align*}
    t_{\text{obs}} &= \frac{\bar{X} - \mu_0}{S/\sqrt{n}} \\
    &= \frac{8.5 - 10}{3.2/\sqrt{20}} \\
    &= \frac{-1.5}{3.2/4.472} \\
    &= \frac{-1.5}{0.715} = -2.098
    \end{align*}
    
    \item \textbf{Degrés de liberté} : $\nu = n - 1 = 20 - 1 = 19$
    
    \item \textbf{Valeurs critiques} (test bilatéral, $\alpha = 0.05$) :
    \begin{align*}
    t_{\alpha/2}(19) = t_{0.025}(19) = 2.093
    \end{align*}
    
    \item \textbf{Décision} :
    \begin{align*}
    |t_{\text{obs}}| = 2.098 > 2.093 = t_{0.025}(19)
    \end{align*}
    
    On \textbf{rejette} $H_0$ au seuil de 5\%. Les données ne confirment pas l'affirmation du laboratoire.
    
    \item \textbf{Calcul de la p-value} :
    \begin{align*}
    p &= 2 \times P(t(19) > |t_{\text{obs}}|) \\
    &= 2 \times P(t(19) > 2.098)
    \end{align*}
    
    Dans la table de Student avec $\nu = 19$ :
    \begin{itemize}
        \item $t_{0.025}(19) = 2.093$
        \item $t_{0.02}(19) = 2.205$
    \end{itemize}
    
    On a $2.093 < 2.098 < 2.205$, donc :
    \begin{align*}
    0.02 < \frac{p}{2} < 0.025 \Rightarrow 0.04 < p < 0.05
    \end{align*}
    
    La p-value est légèrement inférieure à 0.05, ce qui confirme le rejet de $H_0$.
    
    \item \textbf{Interprétation} : La réduction moyenne observée (8.5 mmHg) est significativement différente de 10 mmHg. Le médicament ne produit pas la réduction annoncée.
\end{enumerate}

\begin{attentionbox}
\textbf{Attention} : Quand $n < 30$ et la variance est inconnue, il est crucial d'utiliser la loi de Student et non la loi normale. L'utilisation de la loi normale dans ce cas donnerait des résultats incorrects.
\end{attentionbox}

\newpage

\subsection{Exercices : Test Unilatéral sur la Moyenne}

\subsubsection{Exercice 7 : Test unilatéral à droite avec variance connue}

Un restaurant prétend que ses plats contiennent au moins 250 g de viande. Un inspecteur prélève un échantillon de 40 plats et trouve une moyenne de 248 g. On suppose que le poids suit une loi normale avec un écart-type de 8 g.

\textbf{Question} : Peut-on conclure que le restaurant ne respecte pas sa promesse au seuil de 5\% ?

\subsubsection*{Solution de l'Exercice 7}

\begin{enumerate}
    \item \textbf{Formulation des hypothèses} :
    \begin{align*}
    H_0 &: \mu = 250 \quad \text{(ou } \mu \geq 250\text{)} \\
    H_1 &: \mu < 250 \quad \text{(test unilatéral à gauche)}
    \end{align*}
    
    \textit{Note} : On teste si la moyenne est inférieure à 250, ce qui correspond à un test unilatéral à gauche.
    
    \item \textbf{Données} :
    \begin{itemize}
        \item $n = 40$ (grand échantillon)
        \item $\bar{x} = 248$ g
        \item $\sigma = 8$ g (connu)
        \item $\mu_0 = 250$ g
        \item $\alpha = 0.05$
    \end{itemize}
    
    \item \textbf{Choix de la distribution} : Variance connue $\Rightarrow$ Loi normale $\N(0,1)$
    
    \item \textbf{Calcul de la statistique de test} :
    \begin{align*}
    Z_{\text{obs}} &= \frac{\bar{X} - \mu_0}{\sigma/\sqrt{n}} \\
    &= \frac{248 - 250}{8/\sqrt{40}} \\
    &= \frac{-2}{8/6.325} \\
    &= \frac{-2}{1.265} = -1.581
    \end{align*}
    
    \item \textbf{Valeur critique} (test unilatéral à gauche, $\alpha = 0.05$) :
    \begin{align*}
    z_{1-\alpha} = z_{0.95} = -z_{0.05} = -1.645
    \end{align*}
    
    Dans la table normale : $\Phi(1.645) = 0.95$, donc $P(Z < -1.645) = 0.05$
    
    \item \textbf{Décision} :
    \begin{align*}
    Z_{\text{obs}} = -1.581 > -1.645 = z_{0.95}
    \end{align*}
    
    On \textbf{accepte} $H_0$ au seuil de 5\%. On ne peut pas conclure que le restaurant ne respecte pas sa promesse.
    
    \item \textbf{Calcul de la p-value} :
    \begin{align*}
    p &= P(Z < Z_{\text{obs}}) \\
    &= P(Z < -1.581) \\
    &= \Phi(-1.581) \\
    &= 1 - \Phi(1.581) \\
    &= 1 - 0.9430 = 0.057
    \end{align*}
    
    $p = 0.057 > 0.05 = \alpha$, ce qui confirme qu'on accepte $H_0$.
    
    \item \textbf{Interprétation} : Bien que la moyenne observée (248 g) soit inférieure à 250 g, la différence n'est pas statistiquement significative au seuil de 5\%. On ne peut pas conclure que le restaurant ne respecte pas sa promesse.
\end{enumerate}

\begin{rappelbox}
\textbf{Rappel} : Pour un test unilatéral à gauche ($H_1 : \mu < \mu_0$), on rejette $H_0$ si $Z_{\text{obs}} < z_{1-\alpha}$ (ou $Z_{\text{obs}} < -z_{\alpha}$). Une valeur négative de $Z_{\text{obs}}$ correspond à une moyenne observée inférieure à $\mu_0$.
\end{rappelbox}

\newpage

\subsubsection{Exercice 8 : Test unilatéral à droite avec variance inconnue (loi de Student)}

Un entraîneur prétend que ses athlètes peuvent courir 100 mètres en moins de 12 secondes en moyenne. On teste 15 athlètes et on obtient un temps moyen de 11.8 secondes avec un écart-type de 0.6 secondes. On suppose que les temps suivent une loi normale.

\textbf{Question} : Les données confirment-elles l'affirmation de l'entraîneur au seuil de 1\% ?

\subsubsection*{Solution de l'Exercice 8}

\begin{enumerate}
    \item \textbf{Formulation des hypothèses} :
    \begin{align*}
    H_0 &: \mu = 12 \quad \text{(ou } \mu \geq 12\text{)} \\
    H_1 &: \mu < 12 \quad \text{(test unilatéral à gauche)}
    \end{align*}
    
    \textit{Note} : "Moins de 12 secondes" signifie $\mu < 12$, donc test unilatéral à gauche.
    
    \item \textbf{Données} :
    \begin{itemize}
        \item $n = 15$ (petit échantillon)
        \item $\bar{x} = 11.8$ secondes
        \item $s = 0.6$ secondes
        \item $\mu_0 = 12$ secondes
        \item $\alpha = 0.01$
    \end{itemize}
    
    \item \textbf{Choix de la distribution} : Variance inconnue + $n < 30$ $\Rightarrow$ Loi de Student $t(14)$
    
    \item \textbf{Calcul de la statistique de test} :
    \begin{align*}
    t_{\text{obs}} &= \frac{\bar{X} - \mu_0}{S/\sqrt{n}} \\
    &= \frac{11.8 - 12}{0.6/\sqrt{15}} \\
    &= \frac{-0.2}{0.6/3.873} \\
    &= \frac{-0.2}{0.155} = -1.290
    \end{align*}
    
    \item \textbf{Degrés de liberté} : $\nu = n - 1 = 15 - 1 = 14$
    
    \item \textbf{Valeur critique} (test unilatéral à gauche, $\alpha = 0.01$) :
    \begin{align*}
    t_{1-\alpha}(14) = t_{0.99}(14) = -t_{0.01}(14) = -2.624
    \end{align*}
    
    Dans la table de Student : $t_{0.01}(14) = 2.624$, donc $P(t(14) < -2.624) = 0.01$
    
    \item \textbf{Décision} :
    \begin{align*}
    t_{\text{obs}} = -1.290 > -2.624 = t_{0.99}(14)
    \end{align*}
    
    On \textbf{accepte} $H_0$ au seuil de 1\%. Les données ne confirment pas l'affirmation de l'entraîneur.
    
    \item \textbf{Calcul de la p-value} :
    \begin{align*}
    p &= P(t(14) < t_{\text{obs}}) \\
    &= P(t(14) < -1.290)
    \end{align*}
    
    Dans la table de Student avec $\nu = 14$ :
    \begin{itemize}
        \item $t_{0.10}(14) = 1.345$
        \item $t_{0.15}(14) = 1.076$
    \end{itemize}
    
    On a $1.076 < 1.290 < 1.345$, donc :
    \begin{align*}
    0.10 < P(t(14) > 1.290) < 0.15
    \end{align*}
    
    Par symétrie : $0.10 < p < 0.15$
    
    La p-value est bien supérieure à 0.01, ce qui confirme qu'on accepte $H_0$.
    
    \item \textbf{Interprétation} : Bien que le temps moyen observé (11.8 s) soit inférieur à 12 s, la différence n'est pas statistiquement significative au seuil de 1\%. On ne peut pas conclure que les athlètes courent en moyenne en moins de 12 secondes.
\end{enumerate}

\begin{astucebox}
\textbf{Astuce} : Pour un test unilatéral avec la loi de Student, utilisez la symétrie de la distribution : $P(t(\nu) < -a) = P(t(\nu) > a)$ pour tout $a > 0$.
\end{astucebox}

\newpage

\section{Résumé et Tableau Récapitulatif}

\subsection{Tableau de choix de la distribution}

\begin{table}[h]
\centering
\begin{tabular}{|l|l|l|}
\hline
\textbf{Condition} & \textbf{Distribution} & \textbf{Statistique de test} \\
\hline
Variance connue & Normale $\N(0,1)$ & $Z = \frac{\bar{X} - \mu_0}{\sigma/\sqrt{n}}$ \\
\hline
Variance inconnue, $n < 30$ & Student $t(n-1)$ & $T = \frac{\bar{X} - \mu_0}{S/\sqrt{n}}$ \\
\hline
Variance inconnue, $n \geq 30$ & Normale $\N(0,1)$ (approx.) & $Z = \frac{\bar{X} - \mu_0}{S/\sqrt{n}}$ \\
\hline
Test sur la variance & Khi-deux $\chideux(n-1)$ & $\chideux = \frac{(n-1)S^2}{\sigma_0^2}$ \\
\hline
\end{tabular}
\caption{Choix de la distribution selon les conditions}
\end{table}

\subsection{Tableau des valeurs critiques}

\begin{table}[h]
\centering
\begin{tabular}{|l|l|l|}
\hline
\textbf{Type de test} & \textbf{Loi normale} & \textbf{Loi de Student} \\
\hline
Bilatéral, $\alpha = 0.05$ & $\pm 1.96$ & $\pm t_{0.025}(\nu)$ \\
\hline
Bilatéral, $\alpha = 0.01$ & $\pm 2.576$ & $\pm t_{0.005}(\nu)$ \\
\hline
Unilatéral, $\alpha = 0.05$ & $\pm 1.645$ & $\pm t_{0.05}(\nu)$ \\
\hline
Unilatéral, $\alpha = 0.01$ & $\pm 2.326$ & $\pm t_{0.01}(\nu)$ \\
\hline
\end{tabular}
\caption{Valeurs critiques courantes}
\end{table}

\subsection{Checklist pour résoudre un exercice}

\begin{enumerate}
    \item[$\square$] Lire attentivement l'énoncé
    \item[$\square$] Identifier le paramètre à tester ($\mu$, $\sigma^2$, etc.)
    \item[$\square$] Vérifier les conditions (variance connue/inconnue, taille de l'échantillon)
    \item[$\square$] Choisir la distribution appropriée
    \item[$\square$] Formuler $H_0$ et $H_1$ correctement
    \item[$\square$] Calculer la statistique de test
    \item[$\square$] Trouver la valeur critique dans la table appropriée
    \item[$\square$] Comparer et prendre une décision
    \item[$\square$] Calculer la p-value (optionnel mais recommandé)
    \item[$\square$] Interpréter les résultats dans le contexte
\end{enumerate}

\newpage

\section{Astuces et Conseils Pratiques}

\subsection{Erreurs courantes à éviter}

\begin{enumerate}
    \item \textbf{Confondre unilatéral et bilatéral} :
    \begin{itemize}
        \item Unilatéral : "supérieur à", "inférieur à", "au moins", "au plus"
        \item Bilatéral : "égal à", "différent de", "identique"
    \end{itemize}
    
    \item \textbf{Utiliser la maublvaise distribution} :
    \begin{itemize}
        \item Ne pas utiliser la loi normale quand la variance est inconnue et $n < 30$
        \item Ne pas oublier de multiplier par 2 pour les tests bilatéraux
    \end{itemize}
    
    \item \textbf{Erreurs dans la lecture des tables} :
    \begin{itemize}
        \item Confondre $\alpha$ et $1-\alpha$ dans les tables
        \item Oublier les valeurs absolues pour les tests bilatéraux
        \item Ne pas vérifier les degrés de liberté
    \end{itemize}
    
    \item \textbf{Interprétation incorrecte} :
    \begin{itemize}
        \item "Accepter $H_0$" ne signifie pas "$H_0$ est vraie", mais "on n'a pas assez de preuves pour la rejeter"
        \item Une p-value élevée ne signifie pas que $H_0$ est vraie
    \end{itemize}
\end{enumerate}

\subsection{Conseils pour les examens}

\begin{astucebox}
\textbf{Astuces pour réussir} :
\begin{enumerate}
    \item \textbf{Organisez votre travail} : Suivez toujours les étapes dans l'ordre
    \item \textbf{Notez les formules} : Écrivez d'abord les formules avant de calculer
    \item \textbf{Vérifiez vos calculs} : Un calcul incorrect peut tout fausser
    \item \textbf{Interprétez toujours} : Ne vous contentez pas de rejeter/accepter, expliquez ce que cela signifie
    \item \textbf{Utilisez la p-value} : Elle donne plus d'information que juste rejeter/accepter
    \item \textbf{Entraînez-vous avec les tables} : Familiarisez-vous avec la lecture des tables avant l'examen
\end{enumerate}
\end{astucebox}

\subsection{Formules importantes à retenir}

\begin{align*}
\text{Statistique Z (variance connue)} &: Z = \frac{\bar{X} - \mu_0}{\sigma/\sqrt{n}} \\
\text{Statistique T (variance inconnue)} &: T = \frac{\bar{X} - \mu_0}{S/\sqrt{n}} \\
\text{Statistique } \chideux \text{ (variance)} &: \chideux = \frac{(n-1)S^2}{\sigma_0^2} \\
\text{P-value bilatérale} &: p = 2 \times P(\text{stat} > |\text{obs}|) \\
\text{P-value unilatérale} &: p = P(\text{stat} > \text{obs}) \text{ ou } P(\text{stat} < \text{obs})
\end{align*}

\vspace{2cm}

\begin{center}
\textbf{Fin du document}
\end{center}

\end{document}

